\documentclass[12pt]{beamer}

\usepackage[language=british,
            sortlocale=en-GB,
            backend=biber,
            dashed=false,
            citestyle=authoryear-comp,
            bibstyle=authoryear,
            doi=true,
            useprefix=true,
            sorting=ynt,
            maxbibnames=9,
            maxcitenames=2]{biblatex}

\usepackage{filecontents}

\begin{filecontents*}{refs.bib}

@phdthesis{Bartos1999,
	Author = {Bartos, Huba},
	Location = {Budapest},
	Title = {{Morfoszintaxis {\'e}s interpret{\'a}ci{\'o}: A magyar inflexi{\'o}s jelens{\'e}gek szintaktikai h{\'a}ttere [Morphosyntax and interpretation: the syntactic background of Hungarian inflection]}},
	Year = {1999}}

@article{Coppock2013P,
	Author = {Coppock, Elizabeth},
	Journaltitle = {Natural Language Semantics},
	Pages = {345-371},
	Title = {{A semantic solution to the problem of Hungarian object agreement}},
	Volume = {21},
	Year = {2013}}

@article{Dekany2015,
	Author = {D{\'e}k{\'a}ny, {\'E}va},
	Doi = {10.1007/s11049-014-9278-0},
	Journaltitle = {Natural Language \& Linguistic Theory},
	Title = {The syntax of anaphoric possessives in Hungarian},
	Year = {2015}}

@incollection{denDikken1999,
	Author = {den Dikken, Marcel},
	Booktitle = {Crossing Boundaries: {Advances} in the Theory of Central and {Eastern European} Languages},
	Editor = {Kenesei, Istv{\'a}n},
	Location = {Amsterdam/Philadelphia},
	Pages = {137-178},
	Publisher = {John Benjamins},
	Title = {On the Structural Representation of Possession and Agreement},
	Year = {1999}}

@incollection{Szabolcsi1994,
	Author = {Szabolcsi, Anna},
	Crossref = {KieKis1994},
	Pages = {179-274},
	Title = {The Noun Phrase}}

@collection{KieKis1994,
	Editor = {Kiefer, Ferenc and {\'E}. Kiss, Katalin},
	Location = {New York},
	Publisher = {Academic Press},
	Series = {Syntax and Semantics},
	Title = {The Syntactic Structure of {Hungarian}},
	Volume = {27},
	Year = {1994}}

\end{filecontents*}

\addbibresource{refs.bib}

%\addbibresource{legenda.bib}

\renewcommand\nameyeardelim{\space}
\renewcommand\postnotedelim{\addcolon\space}
\DeclareFieldFormat{pages}{#1}
\DeclareFieldFormat{postnote}{#1} % no postnote prefix in "normal" citation commands
\AtEveryBibitem{%
    \clearfield{labelmonth}%
    \clearfield{issn}%
}

\usepackage{xpatch}

\xpatchbibmacro{name:andothers}{%
  \bibstring{andothers}%
}{%
  \bibstring[\emph]{andothers}%
}{}{}

% No In: for journals

\renewbibmacro{in:}{\ifentrytype{article}{}{\printtext{\bibstring{in}\intitlepunct}}}

% Lower case for prefix in bibliography
% \AtBeginBibliograph{...} is important, otherwise it breaks \Textcite{...}

\AtBeginBibliography{%
\renewbibmacro*{name:last-first}[4]{%
  \ifuseprefix
    {\usebibmacro{name:delim}{#3#1}%
     \usebibmacro{name:hook}{#3#1}%
     \ifblank{#3}{}{%
       %\ifcapital% comment out
       %  {\mkbibnameprefix{\MakeCapital{#3}}\isdot}% and get rid of this
         {\mkbibnameprefix{#3}\isdot}%
       \ifpunctmark{'}{}{\bibnamedelimc}}%
     \mkbibnamelast{#1}\isdot
     \ifblank{#4}{}{\bibnamedelimd\mkbibnameaffix{#4}\isdot}%
     \ifblank{#2}{}{\revsdnamepunct\bibnamedelimd\mkbibnamefirst{#2}\isdot}}
    {\usebibmacro{name:delim}{#1}%
     \usebibmacro{name:hook}{#1}%
     \mkbibnamelast{#1}\isdot
     \ifblank{#4}{}{\bibnamedelimd\mkbibnameaffix{#4}\isdot}%
     \ifblank{#2#3}{}{\revsdnamepunct}%
     \ifblank{#2}{}{\bibnamedelimd\mkbibnamefirst{#2}\isdot}%
     \ifblank{#3}{}{\bibnamedelimd\mkbibnameprefix{#3}\isdot}}}

\renewbibmacro*{name:last}[4]{%
  \ifuseprefix
    {\usebibmacro{name:delim}{#3#1}%
     \usebibmacro{name:hook}{#3#1}%
     \ifblank{#3}
       {}
       {%\ifcapital%<-- commented this out
        %  {\mkbibnameprefix{\MakeCapital{#3}}\isdot}% and this
          {\mkbibnameprefix{#3}\isdot}%
        \ifpunctmark{'}{}{\bibnamedelimc}}}
    {\usebibmacro{name:delim}{#1}%
     \usebibmacro{name:hook}{#1}}%
  \mkbibnamelast{#1}\isdot}%
}

% : as subtitlepunct

\renewcommand{\subtitlepunct}{\addcolon\addspace}

% Sentence case in certain fields

\DeclareFieldFormat{sentencecase}{\MakeSentenceCase{#1}}

\renewbibmacro*{title}{%
\ifthenelse{\iffieldundef{title}\AND\iffieldundef{subtitle}}
{}
{\ifthenelse{\ifentrytype{article}\OR\ifentrytype{inbook}%
  \OR\ifentrytype{incollection}\OR\ifentrytype{inproceedings}%
  \OR\ifentrytype{phdthesis}\OR\ifentrytype{book}}
  {\printtext[title]{%
    \printfield[sentencecase]{title}%
    \setunit{\subtitlepunct}%
    \printfield[sentencecase]{subtitle}}}%
  {\printtext[title]{%
    \printfield[titlecase]{title}%
    \setunit{\subtitlepunct}%
    \printfield[titlecase]{subtitle}}}%
 \newunit}%
\printfield{titleaddon}}

%%% no navigation icons

\setbeamertemplate{navigation symbols}{}

%%% no bibliography icons?

\setbeamertemplate{bibliography item}{}

%% stepwise uncovering

\beamerdefaultoverlayspecification{<+->}

%%% other colors

\definecolor{camcosdark}{RGB}{0,145,193}
\definecolor{camcosblue}{RGB}{0,163,218}
\setbeamercolor*{section in head/foot}{parent=palette quaternary,fg=white,bg=camcosblue}%!90!black}
\setbeamercolor*{subsection in head/foot}{parent=palette secondary,fg=white,bg=camcosblue}
\setbeamercolor{structure}{bg=white,fg=black}
\setbeamercolor*{palette primary}{use=structure,bg=white,fg=black}
\setbeamercolor*{palette secondary}{use=structure,fg=white,bg=camcosblue}%!90!black}
\setbeamercolor*{palette tertiary}{use=structure,fg=white,bg=camcosblue}%!90!black}
\setbeamercolor*{palette quaternary}{use=structure,fg=white,bg=camcosblue}%!90!black}

\setbeamercolor*{alerted text}{fg=camcosblue}%!90!black}
\setbeamercolor*{part in toc}{fg=camcosblue}
\setbeamercolor*{section in toc}{fg=camcosblue}
\setbeamercolor*{item}{fg=camcosblue}

\setbeamertemplate{headline}
{%
\leavevmode
\hbox{\begin{beamercolorbox}[wd=1.00\paperwidth, ht=3.0ex, dp=1.5ex, right, rightskip=.5cm]{section in head/foot}%
        \usebeamerfont{title in head/foot}\scriptsize{\insertsection{}}
    \end{beamercolorbox}

    \begin{beamercolorbox}[wd=.00\paperwidth, ht=2.5ex, dp=1ex, left, leftskip=.5cm]{subsection in head/foot}%
        \usebeamerfont{title in head/foot}%\insertsubsection
    \end{beamercolorbox}
    }
}

\setbeamertemplate{footline}
{%
  \leavevmode%

  \begin{beamercolorbox}[wd=1\paperwidth,ht=2.5ex,dp=1.5ex,leftskip=.3cm plus1fill,rightskip=.3cm]{}%
    \usebeamerfont{title in head/foot}\insertframenumber/\inserttotalframenumber%
  \end{beamercolorbox}%
}

\setbeamertemplate{partpage}
{%
  \begin{centering}
    {\usebeamerfont{part name}\usebeamercolor[bg]{part name}\partname~\insertpartnumber}
    \begin{beamercolorbox}[sep=16pt,center]{part title}
      \usebeamerfont{part title}\insertpart\par
    \end{beamercolorbox}
  \end{centering}
}

\defbeamertemplate*{title page}{customised}[1][]
{%
    \usebeamerfont{title}\alert{\inserttitle}\par
    \usebeamerfont{subtitle}{\large\insertsubtitle}\par
    \bigskip
    \hrule
    \usebeamerfont{author}\insertauthor\par
    \usebeamerfont{institute}\small\insertinstitute\par
    \bigskip
    \hrule
    \usebeamerfont{date}{\normalsize\insertdate}\par
    \usebeamercolor[fg]{titlegraphic}\inserttitlegraphic{}
}

%%%

\setbeamertemplate{caption}[numbered]
\setbeamertemplate{caption label separator}{:}
\setbeamercolor{caption name}{fg=normal text.fg}

\usepackage{amssymb,amsmath}
\usepackage{mathtools}
\usepackage{fixltx2e} % provides \textsubscript

\usepackage{mathspec,xunicode,fontspec,xltxtra}
\defaultfontfeatures{Mapping=tex-text,Scale=MatchLowercase}
\setmainfont[BoldFont = Source Sans Pro Semibold, BoldItalicFont = Source Sans Pro Semibold Italic]{Source Sans Pro}
\setsansfont[BoldFont = Source Sans Pro Semibold, BoldItalicFont = Source Sans Pro Semibold Italic]{Source Sans Pro}
\setmonofont[Scale=0.8]{Source Code Pro}
\setmathfont(Digits,Latin,Greek)[Lowercase=Regular,Uppercase=Regular]{Clear Sans}

\usepackage{microtype}

% linguistics
\usepackage{tikz-qtree}
\usetikzlibrary{positioning,decorations.pathmorphing,decorations,trees,fit}
\tikzstyle{arrow} = [thin, <-, below left, bend right, font = \scriptsize]
\tikzset{>=latex}
\tikzset{every tree node/.style={align=center,anchor=north}}
\tikzset{execute at begin node=\vphantom{A}}
\newcommand*\circled[1]{\tikz[baseline=(char.base)]{
    \node[shape=circle, draw, inner sep=2pt] (char) {#1};}}

\usepackage{pgfplots}

\usepackage{expex}
\lingset{aboveglftskip=-1ex,
        belowglpreambleskip=-1ex,
        belowexskip=-1ex}

\resetcountonoverlays{excnt}

\usepackage{url}

\usepackage{graphicx}
\makeatletter
\def\maxwidth{\ifdim\Gin@nat@width>\linewidth\linewidth\else\Gin@nat@width\fi}
\def\maxheight{\ifdim\Gin@nat@height>\textheight0.8\textheight\else\Gin@nat@height\fi}
\makeatother
% Scale images if necessary, so that they will not overflow the page
% margins by default, and it is still possible to overwrite the defaults
% using explicit options in \includegraphics[width, height, ...]{}
\setkeys{Gin}{width=\maxwidth,height=\maxheight,keepaspectratio}

% Comment these out if you don't want a slide with just the
% part/section/subsection/subsubsection title:
\AtBeginPart{%
  \let\insertpartnumber\relax
  \let\partname\relax
  \frame{\partpage}
}
\AtBeginSection{%
  \let\insertsectionnumber\relax
  \let\sectionname\relax
  \frame{\sectionpage}
}
\AtBeginSubsection{%
  \let\insertsubsectionnumber\relax
  \let\subsectionname\relax
  \frame{\subsectionpage}
}

\usepackage[normalem]{ulem}
% avoid problems with \sout in headers with hyperref:
\pdfstringdefDisableCommands{\renewcommand{\sout}{}}

\setlength{\parindent}{0pt}
\setlength{\parskip}{6pt plus 2pt minus 1pt}
\setlength{\emergencystretch}{3em}  % prevent overfull lines

\setcounter{secnumdepth}{1}

\title{Agreement with possessed direct objects in Hungarian}

\subtitle{An experimental approach}

\author{\textsuperscript{a}András Bárány, \textsuperscript{b}Ádám Szalontai}

\date{SinFonIJA 8, Ljubljana, 25 September 2015}

\institute{\textsuperscript{a}ReCoS/DTAL, University of Cambridge,\\
        \textsuperscript{b}Research Institute for Linguistics,
        Hungarian Academy of Sciences}

\titlegraphic{%
    \begin{tikzpicture}[overlay]
        \node at (8.2, -2.5) {\includegraphics[height=2.5cm]{erc.eps}};
    \end{tikzpicture}
}

\begin{document}

\frame{%
\thispagestyle{empty} % no page number here
  \titlepage{}
}

\begin{frame}{Overview}
\tableofcontents[hideallsubsections]
\end{frame}

\section{Intro: Object agreement and possessed NPs}\label{case-and-case}

\begin{frame}{Object agreement in Hungarian}

\pex
    \a
        \begingl
            \gla    Mari lát egy kutyá-t.//
            \glb    M.      see.\textsc{3sg} a   dog-\textsc{acc}//
            \glft   ‘Mari sees a dog.’//
        \endgl
    \pause
    \a
        \begingl
            \gla    Mari lát-ja a kutyá-t.//
            \glb    M.   see-\textsc{\alert{3sg.obj}} the dog-\textsc{acc}//
            \glft   ‘Mari sees the dog.’//
        \endgl
\xe

    \begin{itemize}

        \item Roughly sensitive to definiteness of the DO

        \item Possessed direct objects are an exception

        \item They can be indefinite yet require agreement

        \item Standard Hungarian: all possessed DOs agree

        \item But: dialectal variation?

    \end{itemize}

\end{frame}

\begin{frame}{This presentation}

    \begin{block}{Questions}

        \begin{itemize}

        \item What kind of variation do we find?

        \item What kinds of factors determine variation?

        \item What determines object agreement?

        \end{itemize}

    \end{block}

    \begin{block}{Methods and results}

        \begin{itemize}

            \item Series of online surveys

            \item Acceptability judgments with and without contexts

            \item Forced choice tests with and without contexts

            \item Speakers treat all possessed DOs alike: \alert{syntactic} trigger?

        \end{itemize}
    \end{block}

\end{frame}

\section{The data}

\begin{frame}{Object agreement}

    \begin{itemize}

        \item Definiteness often a good predictor; not with possessed DOs

        \item Indefinite possessed DOs require agreement

    \end{itemize}

\onslide<3->{%
\pex
    \a
        \begingl
            \gla    Mari {lát /} \ljudge*lát-ja egy kutyá-t.//
            \glb    M.   see.\textsc{3sg} see-\textsc{\alert{3sg.obj}} a   dog-\textsc{acc}//
            \glft   ‘Mari sees a dog.’//
        \endgl
    \pause
    \a
        \begingl
            \gla    Mari {lát-ja /} \ljudge\%lát egy kutyá-m-at.//
            \glb    M.   see-\textsc{\alert{3sg.obj}} see.\textsc{3sg} the dog-\textsc{\alert{1sg.poss}-acc}//
            \glft   ‘Mari sees a dog of mine.’//
        \endgl
\xe}

\end{frame}

\begin{frame}{Types of possessed NPs}

\pex
    \a
        \begingl\rightcomment{Nominative possessor}
            \gla    Mari (*a) kutyá-ja//
            \glb    Mari.\textsc{nom}  (the) dog-\textsc{3sg.poss}//
            \glft   `Mari's dog'//
        \endgl
    \a
        \begingl\rightcomment{Dative possessor}
            \gla    Mari-nak a kutyá-ja//
            \glb    Mari-\textsc{dat} the dog-\textsc{3sg.poss}//
            \glft   `Mari's dog'//
        \endgl
    \a
        \begingl\rightcomment{Pronominal possessor}
            \gla    az ő kutyá-ja//
            \glb    the s/he.\textsc{nom} dog-\textsc{3sg.poss}//
            \glft   `her/his dog'//
        \endgl
\xe

\begin{itemize}

    \item Distribution suggests different heights in the NP \parencite{denDikken1999,Bartos1999,Dekany2015}

\end{itemize}

\end{frame}

\begin{frame}{Structures: \textsc{nom} possessor}

\ex\rightcomment{\emph{Mari kutyája} `Mari's dog'}
    \begin{tikzpicture}[baseline]
        \Tree   [.{\textcolor<2->{camcosblue}{DP}}
                    \node (specdp) {\textcolor<2->{camcosblue}{\emph{Mari.\textsc{nom}}}};
                    [.{D$'$}
                        D
                        [.PossP
                            \node (t) {\emph{\sout{Mari}}};
                                [.{Poss$'$}
                                    [.NP
                                        \node {\emph{kutyája}\\dog};
                                    ]
                                    [.Poss
                                        \node (poss) {\emph{-ja}\\\textsc{3sg.poss}};
                                    ]
                                ]
                        ]
                    ]
                ]
        \draw [arrow, ->, bend left=55] (t.south) to (specdp.south);
    \end{tikzpicture}
\xe

\end{frame}

\begin{frame}{Structures: \textsc{dat} possessor}

\ex\rightcomment{\emph{Marinak a kutyája} `Mari's dog'}
    \begin{tikzpicture}[baseline]
        \Tree   [.{\textcolor<2->{camcosblue}{DP}}
                    \node (specdp) {\textcolor<2->{camcosblue}{\emph{Marinak-\textsc{dat}}}};
                    [.{\textcolor<2->{camcosblue}{DP}}
                        \node (d) {D\\\emph{a}\\the};
                        [.PossP
                            \node (t) {\emph{\sout{Mari}}};
                                [.{Poss$'$}
                                    [.NP
                                        \node {\emph{kutyája}\\dog};
                                    ]
                                    [.Poss
                                        \node (poss) {\emph{-ja}\\\textsc{3sg.poss}};
                                    ]
                                ]
                        ]
                    ]
                ]
        \draw [arrow, ->, bend left=55] (t.south) to (specdp.south);
    \end{tikzpicture}
\xe

\end{frame}

\begin{frame}{Structures: pronominal possessor}

\ex\rightcomment{\emph{az ő kutyája} `her/his dog'}
    \begin{tikzpicture}[baseline]
        \Tree   [.{\textcolor<2->{camcosblue}{DP}}
                    \node (d) {\textcolor<2->{camcosblue}{D}\\\emph{az}\\\emph{the}};
                    [.PossP
                        \node (t) {\textcolor<2->{camcosblue}{\emph{ő}}\\\textcolor<2->{camcosblue}{s/he}};
                            [.{Poss$'$}
                                [.NP
                                    \node {\emph{kutyája}\\dog};
                                ]
                                [.Poss
                                    \node (poss) {\emph{-ja}\\\textsc{3sg.poss}};
                                ]
                            ]
                    ]
                ]
    \end{tikzpicture}
\xe

\end{frame}

\begin{frame}{Syntactic and semantic properties}

    \begin{block}{Extraction}

        \begin{itemize}

            \item Dative possessors can be extracted \parencite{Szabolcsi1994}

            \item Only extracted possessors in definiteness effect-contexts (\nextx a)

            \item Non-specific reading in (\nextx a)

        \end{itemize}

    \end{block}

\onslide<2->{%
\pex
    \a
        \begingl
            \gla    Mari-nak van kutyá-ja.//
            \glb    Mari-\textsc{dat} is dog-\textsc{3sg.poss}//
            \glft   `Mari has a dog/dogs.'//
        \endgl
    \a
        \begingl
            \gla    \ljudge*Mari van kutyá-ja.//
            \glb    Mari.\textsc{nom} is dog-\textsc{3sg.poss}//
        \endgl
\xe}

\end{frame}

\begin{frame}{What triggers object agreement?}

    \begin{block}{Syntax?}
        \begin{itemize}

            \item DP structure\hfill\parencite{Bartos1999}

        \end{itemize}
    \end{block}

    \begin{block}{Semantics?}
        \begin{itemize}

            \item A feature [\textsc{def}] introduces presupposition\hfill\parencite{Coppock2013P}

        \end{itemize}
    \end{block}

    \begin{block}{Both? Neither?}
        \begin{itemize}

            \item Possessed DOs have DP but they can be non-specific?
            \item How does this influence agreement?

        \end{itemize}
    \end{block}

\end{frame}

\section{Surveys}

\begin{frame}{Acknowledgements}

\begin{center}
    \Large \textbf{Thank you!}
\end{center}

The first author's research is supported by ERC Grant No. 269752
\enquote{Rethinking Comparative Syntax}. Thanks to: Theresa Biberauer, Anders
Holmberg, Georg Höhn, Ian Roberts, Michelle Sheehan, and Jenneke van der Wal.

\end{frame}

\begin{frame}[t, allowframebreaks]

    \frametitle{References}
    \newrefcontext[sorting=nyt]
    \renewcommand*{\bibfont}{\footnotesize}
    \printbibliography[heading=none]

\end{frame}

\end{document}
